\documentclass{article}
% Shared layout. Keep benchmark prose and result-specific plots outside this folder.
\PassOptionsToPackage{numbers}{natbib}
\usepackage[preprint]{neurips_2026}
\usepackage[T1]{fontenc}
\usepackage[hidelinks]{hyperref}
\usepackage{booktabs,amsmath,microtype,pgfplots}
\pgfplotsset{compat=1.17}
\usetikzlibrary{patterns}
\definecolor{okblue}{HTML}{0072B2}
\definecolor{okorange}{HTML}{D55E00}
\pgfplotsset{every axis/.append style={font=\small, tick label style={font=\small}, width=\linewidth, height=5cm, grid=major}}

\newcommand{\PaperTitle}{Hundred Step: Dependency Preservation and Recovery in Long Workflows}
\newcommand{\AuthorOrcid}{0009-0005-9964-6841}
\newcommand{\TaskCount}{24}
\newcommand{\SimCount}{4608}
\newcommand{\SeedCount}{24}
\newcommand{\ControlCount}{10}
\newcommand{\ControlEpisodes}{240}
\newcommand{\OracleCompletion}{100.0 [86.2, 100.0]}
\newcommand{\NoCompletion}{0.0 [0.0, 13.8]}
\newcommand{\NoScore}{0.077 [0.050, 0.107]}
\newcommand{\RetryScore}{0.374 [0.319, 0.417]}
\newcommand{\NoSurvival}{31.7 [24.6, 39.0]}
\newcommand{\RetrySurvival}{72.1 [65.8, 76.9]}
\newcommand{\SimShort}{60.4 [53.5, 67.4]}
\newcommand{\SimLong}{13.2 [8.3, 18.8]}
\newcommand{\SimAggregate}{33.9 [30.9, 37.2]}
\newcommand{\HazardLoss}{66.1 [62.8, 69.3]}
\newcommand{\AbstractText}{Long workflows require agents to retain exact dependencies and recover from faults without skipping work. Hundred Step isolates these behaviors in 24 synthetic workflows with replay-based scoring. A reactive scripted controller completes all tasks: 100.0\% (95\% Wilson interval 86.2 to 100.0). Disabling recovery or allowing transient retries alone prevents completion on every task. We also execute 4608 simulated episodes with matched omission draws and a recovery ablation. The stipulated omission probability is 0.02 per logical step. With recovery enabled, simulated completion is 60.4\% (95\% seed-bootstrap interval 53.5 to 67.4) at 25 steps. Simulated completion is 13.2\% (8.3 to 18.8) at 100 steps. The study validates diagnostic behavior under explicit controls. It provides no learned-model leaderboard, human baseline, or estimate of production reliability.}

\title{\PaperTitle}
\author{Pooria Arab\\Independent researcher\\ORCID: \href{https://orcid.org/\AuthorOrcid}{\AuthorOrcid}}
\begin{document}
\maketitle
\begin{abstract}\AbstractText\end{abstract}
\section{Why this matters}
A useful agent must finish work after its first correct action. Each action can create a dependency for the next action. A failed tool call can require a retry, a rollback, or a stop. These cases become hard to distinguish when a workflow contains many steps. An agent that loses the last valid state can repeat work or continue from an invalid checkpoint. A fluent final message does not establish that the work occurred.\par
Hundred Step measures a narrow form of long-horizon reliability: preserving exact dependencies while completing a sequential workflow with recoverable faults. The environment accepts the next operation only when the agent supplies the last valid token. Faults require a specified recovery action before progress can resume. The score comes from replaying actions against the environment. It does not depend on the agent's claim of success.\par
We test whether controlled policy changes produce the intended distinctions in this measurement setting. We separate successful dependency tracking, transient recovery, rollback recovery, and final completion. We also examine a simulated population with stipulated omissions. That experiment shows how the measurement responds when local reliability changes. It does not estimate how often a language model makes those omissions.\par
A long list of actions is not necessarily a demanding planning problem. Here the next action is usually known. Difficulty lies in maintaining state and following the recovery protocol over time. A small deterministic controller can solve the tasks. That fact is a useful positive control, not evidence of broad intelligence. The benchmark tests whether an agent can preserve a simple invariant throughout an extended interaction.\par
\section{Related work}
WebArena evaluates agents through functional tasks in realistic web environments \citep{webarena}. Its setting includes interface navigation and domain context. Hundred Step removes these sources of variation to isolate ordered execution and recovery. This reduction supports diagnosis but sacrifices ecological validity. A strong score here does not imply successful web use.\par
AgentBench studies reasoning and decision making across multiple interactive environments \citep{agentbench}. Hundred Step provides a more restricted axis of evaluation. Its workflow length and fault schedule are explicit. This permits direct inspection of where a controller stops. It cannot represent the breadth of decisions found in a multi-environment evaluation.\par
The $\tau$-bench framework combines tools, policy constraints, and user interaction \citep{taubench}. It also studies consistency across repeated trials. Our deterministic repeats serve a different purpose: they check whether the experiment produces the same trace. They are not additional independent observations. The simulated seeds introduce variation only through an explicitly defined omission process.\par
BetterBench emphasizes benchmark design, documentation, and repeatability \citep{betterbench}. We follow those concerns by disclosing controls, uncertainty units, and limits. We do not claim a formal BetterBench assessment. The result packet makes replay and recomputation possible, but independent validation remains necessary.\par
The design also follows PayBench's construct-first presentation \citep{paybench}. Its useful general principle is to specify the behavior of interest before selecting a headline score. Hundred Step does not adopt payment scenarios, human preference labels, or matched safe and unsafe pairs. Our paired intervention changes a simulated controller's recovery ability while holding task inputs and omission draws fixed.\par
\section{Experiment design}
\subsection{Observable construct}
The intended construct has three observable parts. Dependency preservation means supplying the exact token from the last successful operation. Recovery means applying the action required by a fault before retrying. Completion means reaching the terminal state after every required operation. We observe these properties through tool calls and replayed state.\par
The workflow state contains a completed-step counter, a current token, and any pending fault. An advance request supplies a step number and a token. The environment rejects a request that skips a step or uses the wrong token. A successful request returns the next token. The initial token is supplied openly, so the first action does not require guessing.\par
Transient faults preserve the valid token and request an exact retry. Rollback faults require an acknowledgement with the current checkpoint and its token. A correct acknowledgement permits the failed advance to be retried. The acknowledgement does not undo earlier successful operations. Thus, rollback here tests compliance with a checkpoint protocol. It does not model compensating database transactions or physical reversal of an action.\par
\subsection{Information boundary}
The reactive controller observes the initial token and previous tool outputs. It does not read the schedule of future faults. After a successful advance, it updates its token and step counter. After a transient fault, it repeats the advance. After a rollback fault, it acknowledges the checkpoint and retries. This controller provides a solvability control without consulting the future schedule.\par
The oracle control has a stronger information advantage. It uses the benchmark's reference sequence, including scheduled recovery actions. The oracle tests the scoring ceiling. Agreement between the oracle and reactive controller establishes that advance knowledge is unnecessary for these public tasks. Neither control is a learned model, and neither supplies a competitive baseline for language understanding.\par
The experiment runs each policy through the same public episode runner and replay grader. A policy cannot obtain credit by supplying a forged final state. The replay process reconstructs progress from actions. We challenge that boundary with empty traces accompanied by fabricated success fields. These checks address a concrete false-pass mechanism. They do not prove that every possible malformed trace is rejected.\par
\subsection{Control policies}
The no-recovery policy follows valid tokens but stops when any fault occurs. The retry-only policy handles transient faults but stops when rollback is required. The stale-token policy advances once, retains the initial token, and stops after its next request fails. These policies isolate different parts of the construct while keeping basic task access fixed.\par
The remaining controls submit empty text, a constant answer, an echoed prompt, a refusal, or a single valid first-tool call. They distinguish execution from fluent or relevant text. The first-tool control also checks whether minimal progress receives partial credit without being confused with completion. These controls are deliberately weak. Their failure is necessary evidence for the measurement, but cannot establish its difficulty for capable agents.\par
\subsection{Study order and stopping rule}
The study plan precedes all collected results in version control. Source inspection and the existing family report inform the design. The plan is not externally registered, and version control alone is not an independent timestamp. The local study contains \ControlEpisodes{} control episodes and \SimCount{} simulated episodes before repeats. Every planned cell is retained. Runner failures halt collection instead of becoming missing observations.\par
Each control and each simulated cell is repeated. Semantic hashes compare actions, outputs, replayed state, and scores. Repeats do not increase the sample size. They test reproducibility of the local experiment. No hosted model is configured or probed, so provider availability and learned-model performance remain unmeasured.\par
\section{Coverage}
\begin{figure}[t]\centering
\begin{tikzpicture}\begin{axis}[ybar,bar width=14pt,xlabel={Required logical steps},ylabel={Scheduled faults},xtick={25,50,75,100},ymin=0,legend style={at={(0.02,0.98)},anchor=north west}]
\addplot[fill=okblue] table[x=length,y=transient,col sep=comma]{data/coverage.generated.csv};
\addplot[fill=white,pattern=north east lines] table[x=length,y=rollback,col sep=comma]{data/coverage.generated.csv};
\legend{Transient,Rollback}
\end{axis}\end{tikzpicture}
\caption{Longer tasks carry more transient faults, while rollback coverage stays fixed. Counts describe the complete public suite, not sampled rates. All workflow families appear at each length.}\label{fig:coverage}
\end{figure}

The public suite crosses workflow labels with lengths. The labels describe ticket handling, invoices, inventory, onboarding, migration, and scheduling. They give a familiar interpretation to sequential work. They do not implement those domains. The same abstract state machine governs every label, so these are not independent domain simulators.\par
\begin{table}[t]\centering
\caption{Complete public coverage. Entries are census counts, not estimated rates. Each task contains both fault classes.}\label{tab:coverage}
\begin{tabular}{rrrr}\toprule Required steps & Tasks & Transient faults & Rollback faults\\\midrule
25 & 6 & 6 & 6 \\
50 & 6 & 12 & 6 \\
75 & 6 & 12 & 6 \\
100 & 6 & 18 & 6 \\\bottomrule
\end{tabular}\end{table}
Table~\ref{tab:coverage} gives the exact allocation. Figure~\ref{fig:coverage} shows that fault burden changes with length. Length therefore does not isolate the effect of additional successful actions alone. Longer workflows also contain more transient faults. A study that wants a pure length effect must hold fault count and placement fixed in a separate design.\par
Faults occur inside the workflow rather than at the first or last operation. Within a task, transient faults precede the final scheduled rollback fault. This ordering has a direct interpretation consequence. A policy that retries transient faults can reach much farther than one that stops at the first fault. Its progress does not establish that it can recover from a rollback.\par
The suite uses generated tasks with deterministic tokens and schedules. This supports exact replay and repeatable counterexamples. It does not make the suite a random sample of real work. Public tasks can also be memorized. New seeds can change tokens and schedules, but this study uses only the public suite. No result is presented as held-out generalization.\par
\section{Failure taxonomy}
\paragraph{Dependency loss.} The agent sends a stale or guessed token. The next operation cannot use that state. The environment returns an error without revealing a replacement token. The correct token must come from prior successful output. This distinguishes state retention from probing an error message for a missing answer.
\paragraph{Sequence violation.} The agent requests an operation before its predecessor completes. A correct token alone does not authorize an arbitrary step. The environment requires both the next step number and the current token. Sequence checks therefore prevent a policy from skipping directly to the terminal operation.
\paragraph{Transient abandonment.} The agent stops after a fault that an exact retry would resolve. The trace can contain substantial valid progress, but the task remains unfinished. This failure is distinct from dependency loss because the agent can still possess the correct state. The no-recovery control instantiates this behavior.
\paragraph{Rollback omission.} The agent continues to request an advance when a checkpoint acknowledgement is required, or stops instead of acknowledging it. Repeating the failed advance is insufficient. The retry-only control isolates the stopping version of this failure. The experiment does not test a controller that loops indefinitely at rollback.
\paragraph{Premature termination.} The agent ends before the terminal state even when its last action succeeds. The simulated omission process instantiates this case before a new logical step. The process represents a stipulated stopping hazard. It is not a calibrated model of distraction, context loss, or provider interruption.
\paragraph{Unsupported completion claim.} A final answer or submitted state declares success without the necessary actions. Replay prevents that claim from replacing evidence of execution. Empty and constant-output controls probe this distinction. A low score on these controls cannot rule out more complex scoring attacks.
These classes need not be mutually exclusive in a real trace. A stale token can lead to abandonment, followed by a false completion claim. For diagnosis, the earliest unresolved cause and final outcome should both be retained. The present controlled policies each target a narrow mechanism. Their behavior supplies interpretable examples rather than an exhaustive distribution of natural agent failures.\par
\section{Metrics and scoring}
Let $L$ be the required workflow length and $c$ the number of completed steps. Define survival as $s=c/L$, capped at the last step before an unresolved fault. Let $r$ be resolved injected faults divided by scheduled injected faults. Let $b$ be correct rollback acknowledgements divided by required rollbacks. Let $f$ indicate terminal success. The benchmark composite is\par
\begin{equation}\label{eq:score}
Q=\begin{cases}0,&c=0,\\
(f+(1-f)s)(0.6s+0.1r+0.05b)+0.25f,&c>0.\par
\end{cases}
\end{equation}
The coefficients define the measurement rule; they are not fitted from this study. For unfinished episodes, the leading survival factor discounts partial progress. The survival contribution then becomes quadratic. Full completion receives an additional terminal reward. These choices favor finishing, but they also encode a judgment about partial work. No user study validates the relative weights.\par
Completion is the primary reported outcome because it has a direct state-based meaning. Composite score and survival explain how incomplete policies differ. Recovery and rollback counts remain available in the traces. A task can stop before reaching a scheduled fault. Its recovery denominator still includes that fault. The resulting rate therefore measures coverage of required recoveries, not recovery conditional on encountering an error.\par
For controls, completion intervals use Wilson's binomial construction. Each policy is evaluated on the full public suite of \TaskCount{} tasks. The interval describes uncertainty under a hypothetical exchangeable-task sampling model. The suite is fixed and shares templates, so the interval does not establish population coverage. We report exact completion counts beside it to prevent that interpretation.\par
Control score and survival intervals use a bootstrap over workflow families. A sampled family carries all its lengths. This keeps within-family observations together. The bootstrap cannot represent task families absent from the suite. With few clusters, its tails are approximate. Identical outcomes can produce a collapsed interval, which indicates a boundary in these observations rather than certainty about future tasks.\par
For the simulated experiment, the unit of uncertainty is a simulated seed. Each seed contributes a mean across the same fixed tasks. Bootstrap draws resample these seed summaries and preserve pairing between recovery conditions. This avoids treating every simulated episode as an independent task sample. The resulting intervals quantify Monte Carlo variation under the stipulated process, conditional on the public suite.\par
\section{Results}
\subsection{Controls separate completion from progress}
\begin{figure}[t]\centering
\begin{tikzpicture}\begin{axis}[ybar,ymin=0,ymax=110,ylabel={Completion (\%)},xtick={0,1,2,3,4,5,6,7,8,9},xticklabels={Oracle,Reactive,No recovery,Retry only,Stale token,Empty,Constant,Echo,Refusal,First tool},x tick label style={rotate=35,anchor=east},height=6cm]
\addplot[fill=okblue,error bars/.cd,y dir=both,y explicit] table[x=index,y=value,y error minus=low,y error plus=high,col sep=comma]{data/controls.generated.csv};
\end{axis}\end{tikzpicture}
\caption{Only the oracle and reactive controls finish. Bars show completion across \TaskCount{} tasks per policy, with Wilson reference intervals. These controls are scripted, not learned models.}\label{fig:controls}
\end{figure}

The oracle and reactive controller each complete all public tasks: \OracleCompletion{} percent, with Wilson reference intervals in brackets. This confirms an attainable ceiling through two distinct information paths. The reactive result also shows that the stated recovery protocol is sufficient. It does not show that a language model will infer or maintain that protocol.\par
\begin{table}[t]\centering\small
\caption{Scripted control outcomes. Completion is percent with Wilson intervals; score is a unit-scale mean with workflow-cluster bootstrap intervals. All intervals are nominal 95\%.}\label{tab:controls}
\begin{tabular}{lrlr}\toprule Policy & Finished & Completion [interval] & Score [interval]\\\midrule
oracle & 24/24 & 100.0 [86.2, 100.0] & 1.000 [1.000, 1.000] \\
reactive & 24/24 & 100.0 [86.2, 100.0] & 1.000 [1.000, 1.000] \\
no-recovery & 0/24 & 0.0 [0.0, 13.8] & 0.077 [0.050, 0.107] \\
retry-only & 0/24 & 0.0 [0.0, 13.8] & 0.374 [0.319, 0.417] \\
stale-token & 0/24 & 0.0 [0.0, 13.8] & 0.00034 [0.00034, 0.00034] \\
empty & 0/24 & 0.0 [0.0, 13.8] & 0.000 [0.000, 0.000] \\
constant & 0/24 & 0.0 [0.0, 13.8] & 0.000 [0.000, 0.000] \\
echo & 0/24 & 0.0 [0.0, 13.8] & 0.000 [0.000, 0.000] \\
refusal & 0/24 & 0.0 [0.0, 13.8] & 0.000 [0.000, 0.000] \\
first-tool & 0/24 & 0.0 [0.0, 13.8] & 0.00034 [0.00034, 0.00034] \\\bottomrule
\end{tabular}\end{table}
The no-recovery and retry-only controls both have completion \NoCompletion{} percent. This is a null contrast on the primary outcome. Their partial progress differs substantially. No-recovery has mean survival \NoSurvival{} percent, while retry-only has \RetrySurvival{} percent. Their composite scores are \NoScore{} and \RetryScore{}, respectively. Brackets show the workflow-cluster intervals.\par
This contrast explains why a single progress score can mislead. Retry-only looks substantially better on partial work, yet still fails every complete workflow. The final scheduled rollback creates a shared barrier. Progress distinguishes how far policies get; completion determines whether they satisfy the task. Both measurements are needed to describe this result without exaggeration.\par
The stale-token and first-tool controls obtain the same small composite score, as Table~\ref{tab:controls} shows. Both preserve only the first successful step. The extra invalid call in stale-token does not create progress. Empty, constant, echo, and refusal controls obtain no progress. Their completion intervals remain visible even at the boundary.\par
The forged-state challenge also yields no accepted completion: \NoCompletion{} percent with the same Wilson reference interval. This result is limited to empty traces carrying fabricated success fields. It establishes that those fields alone cannot bypass replay. It is not a general proof of grader security.\par
\subsection{Extra virtual experiment: simulated population and recovery ablation}
\begin{figure}[t]\centering
\begin{tikzpicture}\begin{axis}[xlabel={Required logical steps},ylabel={Simulated completion (\%)},ymin=0,ymax=110,xtick={25,50,75,100},height=6cm,legend style={at={(0.02,0.03)},anchor=south west}]
\addplot[black,mark=square*,error bars/.cd,y dir=both,y explicit] table[x=length,y=value,y error minus=low,y error plus=high,col sep=comma]{data/sim0.generated.csv};
\addplot[okblue,mark=triangle*,error bars/.cd,y dir=both,y explicit] table[x=length,y=value,y error minus=low,y error plus=high,col sep=comma]{data/sim1.generated.csv};
\addplot[okorange,mark=*,error bars/.cd,y dir=both,y explicit] table[x=length,y=value,y error minus=low,y error plus=high,col sep=comma]{data/sim2.generated.csv};
\addplot[black,dashed,mark=diamond*,error bars/.cd,y dir=both,y explicit] table[x=length,y=value,y error minus=low,y error plus=high,col sep=comma]{data/sim3.generated.csv};
\legend{$q=0$,$q=.005$,$q=.01$,$q=.02$}
\end{axis}\end{tikzpicture}
\caption{Simulated omissions accumulate with workflow length despite correct fault recovery. Each point uses all tasks at that length across \SeedCount{} simulated seeds. Bars show seed-cluster bootstrap intervals. Disabled recovery produces no completions at any setting; the main text reports this null separately.}\label{fig:simulation}
\end{figure}

The simulated population adds an omission probability $q$ before each new logical advance. An omission ends the episode. Otherwise, the controller acts correctly until the next decision. Retries and rollback acknowledgements do not receive extra omission draws. This convention makes the simulated hazard depend on required work, rather than on the number of recovery calls.\par
The simulated conditions use the same random stream for a given task and seed. Increasing $q$ changes the acceptance threshold, so stopping events are nested under the same draws. Recovery-enabled and recovery-disabled conditions also share those draws until their behavior diverges. The intervention changes fault handling, not the task schedule. The disabled condition stops at its first fault even if no omission has occurred.\par
\begin{table}[t]\centering
\caption{Simulated aggregate completion and paired recovery gain. Values are percentages or percentage points, with 95\% seed-cluster bootstrap intervals. Disabled recovery has zero completion and a collapsed zero interval in every cell.}\label{tab:simulation}
\begin{tabular}{rll}\toprule Stipulated $q$ & Recovery-enabled completion & Paired recovery gain\\\midrule
0 & 100.0 [100.0, 100.0] & 100.0 [100.0, 100.0] \\
0.005 & 73.1 [70.5, 76.0] & 73.1 [70.5, 76.0] \\
0.01 & 57.1 [53.6, 60.4] & 57.1 [53.6, 60.4] \\
0.02 & 33.9 [30.9, 37.2] & 33.9 [30.9, 37.2] \\\bottomrule
\end{tabular}\end{table}
With recovery enabled at $q=.02$, simulated completion is \SimShort{} percent for the shortest workflows and \SimLong{} percent for the longest. Across lengths, simulated completion is \SimAggregate{} percent. The paired simulated loss from $q=0$ to $q=.02$ is \HazardLoss{} percentage points. All intervals in this paragraph use seed-cluster bootstrap draws.\par
The simulated recovery-disabled cells remain at the floor throughout the sweep. This null exposes a measurement limit: once every task fails at a scheduled fault, completion cannot expose additional omission damage. Survival can still change, but the completion metric is saturated at failure. A recovery study that uses only these cells cannot rank omission severity through final success.\par
Under the stipulated independent omission process, recovery-enabled completion has expectation $(1-q)^L$. The observed simulated length trend is consistent with that mechanism. The equation is a consequence of the simulation assumptions, not a fitted discovery about agents. Correlated failures, persistent confusion, and recovery overhead could produce different curves. The experiment verifies sensitivity to a known intervention and makes its assumptions inspectable.\par
\section{Roadmap}
The next measurement step is a learned-model study with fixed prompts, budgets, and versioned model identifiers. It should retain the reactive controller as a positive control and the weak policies as negative controls. Provider failures must be distinguished from task failures. A hosted run would also need recorded settings and terms that permit publication of scores. None of those results are supplied here.\par
A second extension should separate horizon length from fault burden. Matched tasks could retain a fixed number of faults while moving their locations. Another condition could hold local error risk constant while changing recovery-call overhead. These interventions would distinguish dependency retention from exposure to additional failure opportunities. They should be specified before collecting comparative model data.\par
A third extension should broaden state and recovery semantics. Branching workflows, stale reads, partial side effects, and true compensating actions can challenge policies that solve this linear protocol. Such additions should keep a transparent reference controller and replayable outcomes. Greater realism alone is not a guarantee of a better measurement. Each new mechanism requires a clear relation to the intended construct.\par
Finally, future releases should use fresh private seeds for final comparisons and disclose them after evaluation. Independent replication should check whether public-task performance transfers to those instances. A version change that alters recovery semantics or scoring requires separate results. Scores from different measurement rules should not be combined into a single trend.\par
\section{Caveats and claim boundaries}
\paragraph{Synthetic domains.} The workflow labels do not instantiate real services. Success shows control over a small protocol, not competence in accounting, identity management, scheduling, or database migration. There are no people, external accounts, or irreversible actions in the evaluation environment.
\paragraph{Public and small coverage.} The fixed suite contains shared templates and a narrow range of failures. Wilson and bootstrap intervals do not correct that selection. A learned agent could exploit public structure. This study neither measures contamination nor establishes a distribution of unseen workflows.
\paragraph{Simulated population.} The simulated seeds represent random omission paths, not people or trained agents. The simulated process assumes independent per-step omissions and otherwise correct behavior. It omits correlated memory failure, tool latency, ambiguity, and strategic recovery. The simulated results must not be used as predicted model scores.
\paragraph{Asymmetric interventions.} Disabling recovery makes every fault-bearing task impossible for that policy. Its failure is expected by construction. The ablation checks causal control within the simulator, but does not estimate the effect of a real product safeguard. Recovery is an ability of the policy here, not an external enforcement layer.
\paragraph{Scoring judgment.} The composite weights and terminal reward reflect design choices. Partial progress can have different value in different applications. We therefore prioritize completion and show components. A user should select a score interpretation before comparing systems, rather than choose the metric that favors a preferred result.
\paragraph{Validation limits.} Semantic repeat agreement supports deterministic reproduction. It does not establish that the task generator or grader matches an independent specification. The public replay path is shared across scoring and verification. Independent implementation, broader trace attacks, and human review remain useful checks against shared mistakes.
\paragraph{Absent evidence.} No free hosted inference is configured or tested in this study. This is not evidence that free inference is unavailable. There is no learned-model leaderboard, human baseline, user survey, production deployment, or independent external registration. The study reports local control evidence and a labelled simulated experiment only.
\section{Citation and availability}
The accompanying packet contains executable experiments, raw traces, generated summaries, and reproducible figures. Code and task data are available at \url{https://github.com/pooriaarab/benchmark-hundred-step}. The published paper is available at \url{https://doi.org/10.5281/zenodo.23752401}; the packet includes structured citation metadata. The paper and generated figures use CC-BY-4.0; benchmark software retains its existing MIT license.\par
\section{Authors and disclosure}
Pooria Arab is an independent researcher and maintains the benchmark. The author therefore has an interest in its usefulness. AI assistance supports the study implementation and manuscript preparation. No independent reviewer or human evaluation is claimed. Local scripted inference incurs no provider charge. The author remains responsible for interpreting the results and correcting errors.\par
\bibliographystyle{plainnat}\bibliography{references}
\end{document}
